\documentclass[a4paper, 14pt]{extarticle}

% Поля
%--------------------------------------
\usepackage{geometry}
\geometry{a4paper,tmargin=2cm,bmargin=2cm,lmargin=3cm,rmargin=1cm}
%--------------------------------------

\usepackage{float}
%Russian-specific packages
%--------------------------------------
\usepackage[T2A]{fontenc}
\usepackage[utf8]{inputenc} 
\usepackage[english, main=russian]{babel}
%--------------------------------------

\usepackage{textcomp}

% Красная строка
%--------------------------------------
\usepackage{indentfirst}               
%--------------------------------------             
\usepackage[most]{tcolorbox}
\tcbuselibrary{listings,breakable}
\newtcblisting{breakablelisting}{
  listing only,
  breakable,
  enhanced,
  frame hidden,
  colback=white,
  listing options={
    mathescape=false,
    texcl=false,
    inputencoding=utf8,
    literate={\$}{{\$}}1
  }
}


%Graphics
%--------------------------------------
\usepackage{graphicx}
\usepackage{wrapfig}
%--------------------------------------
\graphicspath{{./lab5/}}
% Полуторный интервал
%--------------------------------------
\linespread{1.3}                    
%--------------------------------------

%Выравнивание и переносы
%--------------------------------------
% Избавляемся от переполнений
\sloppy
% Запрещаем разрыв страницы после первой строки абзаца
\clubpenalty=10000
% Запрещаем разрыв страницы после последней строки абзаца
\widowpenalty=10000
%--------------------------------------

%Списки
\usepackage{enumitem}

%Подписи
\usepackage{caption} 

%Гиперссылки
\usepackage{hyperref}

\hypersetup {
	unicode=true
}

%Рисунки
%--------------------------------------
\DeclareCaptionLabelSeparator*{emdash}{~--- }
\captionsetup[figure]{labelsep=emdash,font=onehalfspacing,position=bottom}
%--------------------------------------

\usepackage{tempora}

%Листинги
%--------------------------------------
\usepackage{listings}
\lstset{
  basicstyle=\ttfamily\footnotesize, 
  %basicstyle=\footnotesize\AnkaCoder,        % the size of the fonts that are used for the code
  breakatwhitespace=false,         % sets if automatic breaks shoulbd only happen at whitespace
  breaklines=true,                 % sets automatic line breaking
  captionpos=t,                    % sets the caption-position to bottom
  inputencoding=utf8,
  frame=single,                    % adds a frame around the code
  keepspaces=true,                 % keeps spaces in text, useful for keeping indentation of code (possibly needs columns=flexible)
  keywordstyle=\bf,       % keyword style
  numbers=left,                    % where to put the line-numbers; possible values are (none, left, right)
  numbersep=5pt,                   % how far the line-numbers are from the code
  xleftmargin=25pt,
  xrightmargin=25pt,
  showspaces=false,                % show spaces everywhere adding particular underscores; it overrides 'showstringspaces'
  showstringspaces=false,          % underline spaces within strings only
  showtabs=false,                  % show tabs within strings adding particular underscores
  stepnumber=1,                    % the step between two line-numbers. If it's 1, each line will be numbered
  tabsize=2,                       % sets default tabsize to 8 spaces
  title=\lstname                   % show the filename of files included with \lstinputlisting; also try caption instead of title
}
%--------------------------------------

%%% Математические пакеты %%%
%--------------------------------------
\usepackage{amsthm,amsfonts,amsmath,amssymb,amscd}  % Математические дополнения от AMS
\usepackage{mathtools}                              % Добавляет окружение multlined
\usepackage[perpage]{footmisc}
%--------------------------------------

%--------------------------------------
%			НАЧАЛО ДОКУМЕНТА
%--------------------------------------

\begin{document}

%--------------------------------------
%			ТИТУЛЬНЫЙ ЛИСТ
%--------------------------------------
\begin{titlepage}
\thispagestyle{empty}
\newpage


%Шапка титульного листа
%--------------------------------------
\vspace*{-60pt}
\hspace{-65pt}
\begin{minipage}{0.3\textwidth}
\hspace*{-20pt}\centering
\includegraphics[width=\textwidth]{emblem}
\end{minipage}
\begin{minipage}{0.67\textwidth}\small \textbf{
\vspace*{-0.7ex}
\hspace*{-6pt}\centerline{Министерство науки и высшего образования Российской Федерации}
\vspace*{-0.7ex}
\centerline{Федеральное государственное автономное образовательное учреждение }
\vspace*{-0.7ex}
\centerline{высшего образования}
\vspace*{-0.7ex}
\centerline{<<Московский государственный технический университет}
\vspace*{-0.7ex}
\centerline{имени Н.Э. Баумана}
\vspace*{-0.7ex}
\centerline{(национальный исследовательский университет)>>}
\vspace*{-0.7ex}
\centerline{(МГТУ им. Н.Э. Баумана)}}
\end{minipage}
%--------------------------------------

%Полосы
%--------------------------------------
\vspace{-25pt}
\hspace{-35pt}\rule{\textwidth}{2.3pt}

\vspace*{-20.3pt}
\hspace{-35pt}\rule{\textwidth}{0.4pt}
%--------------------------------------

\vspace{1.5ex}
\hspace{-35pt} \noindent \small ФАКУЛЬТЕТ\hspace{80pt} <<Информатика и системы управления>>

\vspace*{-16pt}
\hspace{47pt}\rule{0.83\textwidth}{0.4pt}

\vspace{0.5ex}
\hspace{-35pt} \noindent \small КАФЕДРА\hspace{50pt} <<Теоретическая информатика и компьютерные технологии>>

\vspace*{-16pt}
\hspace{30pt}\rule{0.866\textwidth}{0.4pt}
  
\vspace{11em}

\begin{center}
\Large {\bf Лабораторная работа № 9} \\ 
\large {\bf по курсу <<Методы оптимизации>>} \\
\large <<Метод Ньютона и его модификация>> 
\end{center}\normalsize

\vspace{8em}


\begin{flushright}
  {Студент группы ИУ9-82Б Шемякин В.А. \hspace*{15pt}\\ 
  \vspace{2ex}
  Преподаватель Посевин Д. П.\hspace*{15pt}}
\end{flushright}

\bigskip

\vfill
 

\begin{center}
\textsl{Москва 2026}
\end{center}
\end{titlepage}
%--------------------------------------
%		КОНЕЦ ТИТУЛЬНОГО ЛИСТА
%--------------------------------------

\renewcommand{\ttdefault}{pcr}

\setlength{\tabcolsep}{3pt}
\newpage
\setcounter{page}{2}

\section{Задание}\label{Sect::task}
Необходимо реализовать метод Ньютона, а также модифицированный метод с постоянным Гессианом.
\section{Результаты}\label{Sect::res}

Исходный код программы представлен в листинге 1.

\begin{breakablelisting}
using Plots
using LinearAlgebra
using ForwardDiff

rosenbrock(x) = (1.0 - x[1])^2 + 100.0 * (x[2] - x[1]^2)^2
rastrigin(x) = 20.0 + (x[1]^2 - 10.0 * cos(2π * x[1])) + (x[2]^2 - 10.0 * cos(2π * x[2]))
schwefel(x) = 418.9829 * 2.0 - x[1] * sin(sqrt(abs(x[1]))) - x[2] * sin(sqrt(abs(x[2])))

function find_interval(φ, h=0.1, factor=2.0, max_iter=20)
    a, fa = 0.0, φ(0.0)
    b, fb = h, φ(h)
    if fb > fa
        h = -h
        b, fb = h, φ(h)
        fb > fa && return (-abs(h), abs(h))
    end
    c = b + h
    for _ in 1:max_iter
        fc = φ(c)
        fc > fb && return (min(a, c), max(a, c))
        a, fa, b, fb = b, fb, c, fc
        h *= factor
        c = b + h
    end
    return (a, c)
end

function line_search_1d(f, x, d; tol=1e-8)
    φ(α) = f(x .+ α .* d)
    a, b = find_interval(φ)
    abs(b - a) < 1e-12 && return 0.5 * (a + b)
    φ_gold = (sqrt(5.0) - 1.0) / 2.0
    c = b - φ_gold * (b - a)
    d_pt = a + φ_gold * (b - a)
    fc, fd = φ(c), φ(d_pt)
    while abs(b - a) > tol
        if fc < fd
            b, d_pt, fd = d_pt, c, fc
            c = b - φ_gold * (b - a)
            fc = φ(c)
        else
            a, c, fc = c, d_pt, fd
            d_pt = a + φ_gold * (b - a)
            fd = φ(d_pt)
        end
    end
    return 0.5 * (a + b)
end

function newton_minimize(f, x0; max_iter=100, tol=1e-8)
    x = copy(x0)
    trajectory = [copy(x)]
    for _ in 1:max_iter
        g = ForwardDiff.gradient(f, x)
        H = ForwardDiff.hessian(f, x)
        if norm(g) <= tol
            break
        end
        d = try
            -H \ g
        catch
            -g
        end
        x_new = x .+ d
        if f(x_new) >= f(x)
            α = line_search_1d(f, x, d)
            x_new = x .+ α .* d
        end
        x = x_new
        push!(trajectory, copy(x))
    end
    return x, trajectory
end

function newton_mod_minimize(f, x0; max_iter=100, tol=1e-8)
    x = copy(x0)
    trajectory = [copy(x)]
    H = ForwardDiff.hessian(f, x)
    H_inv = try
        inv(H)
    catch
        n = length(x)
        Matrix{Float64}(I, n, n)
    end
    for _ in 1:max_iter
        g = ForwardDiff.gradient(f, x)
        if norm(g) <= tol
            break
        end
        d = -H_inv * g
        dot(d, g) >= 0 && (d = -g)
        α = line_search_1d(f, x, d)
        x = x .+ α .* d
        push!(trajectory, copy(x))
    end
    return x, trajectory
end

function plot_trajectory(f, trajectory, xlim, ylim, title, filename)
    xs = range(xlim[1], xlim[2], length=200)
    ys = range(ylim[1], ylim[2], length=200)
    Z = [f([x, y]) for y in ys, x in xs]
    plt = contour(xs, ys, Z, levels=20)
    tr = trajectory
    plot!(plt, [p[1] for p in tr], [p[2] for p in tr], color=:red, linewidth=2, label="")
    scatter!(plt, [tr[1][1]], [tr[1][2]], color=:green, markersize=10, label="старт")
    scatter!(plt, [tr[end][1]], [tr[end][2]], color=:blue, markersize=10, label="конец")
    title!(plt, title)
    xlabel!(plt, "x_1")
    ylabel!(plt, "x_2")
    savefig(plt, filename)
end

x0_rastrigin = [0.3, 0.3]
x0_schwefel = [400.0, 400.0]
x0_rosenbrock = [-1.5, 2.0]

lab9_dir = @__DIR__

println("Растригин Ньютон")
x_opt, traj = newton_minimize(rastrigin, x0_rastrigin)
println("  x* = ", x_opt, "\n  f(x*) = ", rastrigin(x_opt), "\n  итераций: ", length(traj) - 1)
rx = (min(0.75, minimum(p[1] for p in traj)), max(1.02, maximum(p[1] for p in traj)))
ry = (min(-0.05, minimum(p[2] for p in traj)), max(0.28, maximum(p[2] for p in traj)))
plot_trajectory(rastrigin, traj, rx, ry, "Растригин Ньютон", joinpath(lab9_dir, "rastrigin_newton.png"))

println("Растригин Ньютон модификация")
x_opt, traj = newton_mod_minimize(rastrigin, x0_rastrigin)
println("  x* = ", x_opt, "\n  f(x*) = ", rastrigin(x_opt), "\n  итераций: ", length(traj) - 1)
rx = (min(0.75, minimum(p[1] for p in traj)), max(1.02, maximum(p[1] for p in traj)))
ry = (min(-0.05, minimum(p[2] for p in traj)), max(0.28, maximum(p[2] for p in traj)))
plot_trajectory(rastrigin, traj, rx, ry, "Растригин Ньютон модификация", joinpath(lab9_dir, "rastrigin_newton_mod.png"))

println("Швефель Ньютон")
x_opt, traj = newton_minimize(schwefel, x0_schwefel)
println("  x* = ", x_opt, "\n  f(x*) = ", schwefel(x_opt), "\n  итераций: ", length(traj) - 1)
plot_trajectory(schwefel, traj, (350, 450), (350, 450), "Швефель Ньютон", joinpath(lab9_dir, "schwefel_newton.png"))

println("Швефель Ньютон модификация")
x_opt, traj = newton_mod_minimize(schwefel, x0_schwefel)
println("  x* = ", x_opt, "\n  f(x*) = ", schwefel(x_opt), "\n  итераций: ", length(traj) - 1)
plot_trajectory(schwefel, traj, (350, 450), (350, 450), "Швефель Ньютон модификация", joinpath(lab9_dir, "schwefel_newton_mod.png"))

println("Розенброк Ньютон")
x_opt, traj = newton_minimize(rosenbrock, x0_rosenbrock)
println("  x* = ", x_opt, "\n  f(x*) = ", rosenbrock(x_opt), "\n  итераций: ", length(traj) - 1)
plot_trajectory(rosenbrock, traj, (-2, 2), (-1, 5), "Розенброк Ньютон", joinpath(lab9_dir, "rosenbrock_newton.png"))

println("Розенброк Ньютон модификация")
x_opt, traj = newton_mod_minimize(rosenbrock, x0_rosenbrock)
println("  x* = ", x_opt, "\n  f(x*) = ", rosenbrock(x_opt), "\n  итераций: ", length(traj) - 1)
plot_trajectory(rosenbrock, traj, (-2, 2), (-1, 5), "Розенброк Ньютон модификация", joinpath(lab9_dir, "rosenbrock_newton_mod.png"))
\end{breakablelisting}

\newpage

Результат представлен на рисунках~\ref{fig:img1}--~\ref{fig:rastrigin}.

\begin{figure}[H]
  \centering
  \includegraphics[width=0.6\textwidth]{lab9/screen1.png}
  \caption{Результат}
  \label{fig:img1}
\end{figure}

\begin{figure}[H]
  \centering
  \includegraphics[width=0.6\textwidth]{lab9/screen2.png}
  \caption{Результат}
  \label{fig:schwefel}
\end{figure}

\begin{figure}[H]
  \centering
  \includegraphics[width=0.6\textwidth]{lab9/screen3.png}
  \caption{Результат}
  \label{fig:rastrigin}
\end{figure}

\end{document}
